\documentclass[10pt,a4paper]{amsart}
\usepackage[margin=19mm]{geometry}
\usepackage{amsmath,amssymb,mathtools}
\usepackage[scr=rsfs,cal=euler]{mathalfa}
\usepackage{xfrac}
\usepackage[unicode,hidelinks,bookmarks=false]{hyperref}
\newcommand{\ie}{i.e.\ }
\newcommand{\cf}{cf.\ }
\newcommand{\resp}{respectively\ }
\newcommand{\Subsection}{\subsection}
\providecommand{\nobreakdash}{\nobreak\hbox{-}}
\setlength{\emergencystretch}{2em}
\multlinegap0pt
\usepackage{amssymb}
\usepackage{xcolor}
\usepackage[normalem]{ulem}
\usepackage[shortlabels]{enumitem}
\usepackage{tikz}
\usepackage{algorithm}\usepackage{algpseudocode}
\usepackage{float}
\newtheorem{lem}{Lemma}
\newcommand{\R}{\mathbb{R}}
\renewcommand{\geq}{\geqslant}
\renewcommand{\leq}{\leqslant}
\title{Incremental Learning in Mirror Flows}
\author{Rapha\"el Berthier, Loucas Pillaud-Vivien}
\date{Source: arXiv:2606.23198v2 (CC BY 4.0)}
\begin{document}
\maketitle

\noindent\emph{Source and licence.} Incremental Learning in Mirror Flows. Version arXiv:2606.23198v2 (CC BY 4.0), distributed by arXiv under the Creative Commons Attribution 4.0 International licence, http://creativecommons.org/licenses/by/4.0/, as recorded in the arXiv licence metadata for this exact version and shown on the abstract page https://arxiv.org/abs/2606.23198v2. Full licence notice: \texttt{licenses/arxiv-2606.23198-CC-BY-4.0.txt}.\par\smallskip

\noindent\textit{Editorial scope.} Counted unit: the article's lem lem:proof-support-function-simplex (source0.tex lines 1542--1549). The article's complete original proof of it is delivered with it (source0.tex lines 1551--1575, one \texttt{proof} environment of 25 lines). No mathematics is written, repaired or re-derived: the statement and the proof are the article's own lines, in the article's own notation, and no retained line is substituted or re-flowed.  No further source-local context is needed: every cross-reference of every retained line resolves inside the two delivered spans.  Nothing was omitted from any retained span, so the package carries no omission record.  No retained line contains a citation, so the wrapper carries no bibliography and no bibliography entry.  No retained line contains a figure environment, an image-inclusion command or drawing code, so no image is delivered and the package declares no asset dependency.  Editorial formatting only, in addition: an AMS \texttt{amsart} preamble that restates the article's own declarations and macros used by the retained lines, namely the environments \texttt{lem} on separate counters, together with the standard packages those lines need.  The retained environments print in this document as Lemma 1 in order of appearance, whereas the article prints the same items with the numbering of its own full document.  The complete native source of the article as distributed in the arXiv e-print for this exact version is bundled unchanged as \texttt{sources/203415/source0.tex}.
\begin{lem}
	\label{lem:proof-support-function-simplex}
	For all $w \in \R^d$, 
	\begin{equation*}
		\partial \sigma_{\Delta_d}(w) = \left\{ \sum_{k \in I(w)} \alpha_k e_k \,\middle|\, \sum_{k \in I(w)} \alpha_k = 1,\, \alpha_k \geq 0 \right\} ,
	\end{equation*}
	where $(e_k)_{1 \leq k \leq d}$ are the canonical vectors.
\end{lem}

\begin{proof}
	In this proof, for any vector $w \in \R^d$, we write $w_{\max} := \max_i w_i = \sigma_{\Delta_d}(w)$. Note that
	\begin{align*}
		p \in \partial \sigma_{\Delta_d}(w) &\iff \forall v \in \R^d, v_{\max} \geq w_{\max} + \langle p, v - w\rangle
	\end{align*}
	We show the lemma by double inclusion.

	First, let $p$ be of the form $\sum_{k \in I(w)} \alpha_k e_k$ where $\sum_{k \in I(w)} \alpha_k = 1$ and $\alpha_k \geq 0$. Let $v \in \R^d$. We have
	%
	\begin{align*}
	\langle p, v - w\rangle &= \sum_{k \in I(w)} \alpha_k (v_k - w_k) = \sum_{k \in I(w)} \alpha_k v_k - w_{\max} \leq v_{\max} - w_{\max}\,,
	\end{align*}
	hence $\partial \sigma_{\Delta_d}(w) \supseteq \left\{ \sum_{k \in I(w)} \alpha_k e_k \,\middle|\, \sum_{k \in I(w)} \alpha_k = 1,\, \alpha_k \geq 0 \right\} $.

	Second, let $p \in \partial \sigma_{\Delta_d}(w)$.
	
	We first show that $p \in \Delta_d$. Let $u \in \R^d$, take $v = w + t u$ for some $t > 0$. We have that 
	%
	\begin{align*}
	\max_{i} \{w_i + t u_i\} &\geq w_{\max} + t\langle p, u\rangle \, ,
	\end{align*}
	Dividing by $t$ and letting $t \to \infty$, as $\max_{i} \{w_i + t u_i\} / t \to \max_i u_i$, we obtain $\langle p, u\rangle \leq \max_i u_i$. Since this holds for all $u \in \R^d$, evaluating the inequality at $u = -e_k$ gives $p \in \R_{\geq 0}^d$, and evaluating it at $u = \pm \mathbf{1}$ gives $p \in \mathcal{S}_1$. Hence $p \in \Delta_d$.

	We now show that $p$ is supported on the active coordinates of $w$. Taking $v = 0$, we obtain $\langle p, w\rangle \geq w_{\max}$. Since $p \in \Delta_d$, we also have $\langle p, w\rangle \leq w_{\max}$. Hence $\langle p, w\rangle = w_{\max}$. Since $p \in \Delta_d$, this can happen only if $p$ is supported on the active coordinates of $w$. This finishes the proof of the second inclusion, and thus of the lemma.
\end{proof}

\end{document}
